\honest{Neural Categories}{Eliezer Yudkowsky}{2008}

Back to bleggs and rubes, with five features: colour, shape, texture, luminance and interior. Suppose I want a network that predicts a blegg's unseen features from its seen ones, and I am a naive designer, unable to derive gradient descent for a simple multilayer network (easier than it sounds), who has read excited popular books saying that neural networks are ``distributed, emergent, and parallel just like the human brain!!''

My first design, Network 1, connects every feature to every other; a twin, Network 1b, sorts humans from Space Monsters with inputs from Aristotle and Plato's Academy. It learns by Hebb's rule, one of the first proposed: features active together get a stronger connection, so seeing blue furred things links blue to fur. Shown something furred, egg-shaped and reddish purple, its activations bounce around, ``in parallel!! and asynchronously!!,'' until it settles on ``glows'' and ``vanadium,'' though no unit stands for the category. I mock this as ``emergent behavior!!'' \nb{The design is the Hopfield network (1982), a standard model of memory, whose abstract speaks of properties that ``emerge as collective properties'' and of ``asynchronous parallel processing.''}

It has problems. Recurrent networks can oscillate, go chaotic or settle slowly, which is bad when you must wait five minutes to settle on ``tiger.'' Evidence is counted twice as activation bounces back and forth: suspecting a glow raises belief in vanadium, which raises belief in the glow. And the connections grow as the square of the number of features. \nb{With the symmetric connections Hebb's rule gives and one-at-a-time updates, Hopfield proved that such a network always settles.}

Network 2 connects each feature only to a central category unit. It answers in one pass, which matters when neurons run at 20Hz, and its connections grow in proportion to the features. \nb{This is the Naive Bayes model, as Yudkowsky wrote three weeks later.} Network 1 can store a pattern Network 2 cannot, such as red objects that never glow, by ``a very strong direct negative connection from color to luminance.'' \nb{On the figure's coding the connection needs to be positive; a negative one would make red objects glow.} That is no special exception, since Network 1 has no blegg unit. But I judge that the extra connections buy little, since real animals are rarely halfway between cat and dog. Some facts neither network can store without extra nodes: sea-blue and spheroid together meaning palladium, but each alone strongly meaning vanadium. Each design assumes something about the world's structure, and reading off those assumptions is what separates adults from babes in machine learning.

Neither design is realistic. ``But it still seems like a fair guess'' that the brain ``is in some sense closer to Network 2,'' because it is ``Fast, cheap, scalable,'' and natural selection likes that. \nb{These are design virtues, not evidence about brains, and the post cites none. In 1994 Murphy and Ross found evidence against the assumption, built into Network 2, that features are independent within a category.}

A brain like that would probably not notice that sea-blue objects never glow, unless twenty of them were shown together and the lights switched off, perhaps forming a subcategory; scattered among a hundred others, they would pass unnoticed. It would decide once that Socrates is human, from his lack of feathers, broad nails, upright walk and Greek, and infer the rest. It would not think to ask how much wearing clothes, as against using language, goes with mortality. Are there biases in classifying things once and for all? ``Of course there are.'' See ``Cultish Countercultishness.'' To be continued. \nb{Murphy and Ross also found that when people are unsure of an object's category, they generally ignore information outside the most likely one.}
